\documentclass[11pt]{article}

\usepackage[margin=0.9in]{geometry}
\usepackage[T1]{fontenc}
\usepackage{lmodern}
\usepackage{amsmath}
\usepackage{xcolor}
\usepackage{array}
\usepackage{booktabs}
\usepackage{longtable}
\usepackage{tabularx}
\usepackage{placeins}
\usepackage{enumitem}
\usepackage{hyperref}
\usepackage{titlesec}
\usepackage{fancyhdr}
\usepackage{graphicx}

\definecolor{prismblue}{HTML}{2457A6}
\definecolor{prismpurple}{HTML}{7562C8}
\definecolor{prismgreen}{HTML}{2E8B57}
\definecolor{darkgray}{HTML}{333B4F}

\hypersetup{
  colorlinks=true,
  linkcolor=prismblue,
  urlcolor=prismpurple,
  citecolor=prismgreen
}

\pagestyle{fancy}
\fancyhf{}
\lhead{Formbot: An Agentic RAG Platform for Form Understanding}
\rhead{Internship Report}
\cfoot{\thepage}
\setlength{\headheight}{14pt}
\renewcommand{\headrulewidth}{0.4pt}

\titleformat{\section}{\Large\bfseries\color{prismblue}}{\thesection}{0.6em}{}
\titleformat{\subsection}{\large\bfseries\color{prismpurple}}{\thesubsection}{0.6em}{}
\titleformat{\subsubsection}{\bfseries\color{darkgray}}{\thesubsubsection}{0.6em}{}

\setlength{\parindent}{0pt}
\setlength{\parskip}{0.55\baselineskip}
\setlist[itemize]{leftmargin=1.4em,itemsep=0.2em,topsep=0.2em}
\setlist[enumerate]{leftmargin=1.6em,itemsep=0.2em,topsep=0.2em}
\renewcommand{\arraystretch}{1.2}

\newcolumntype{Y}{>{\raggedright\arraybackslash}X}

\title{\bfseries\color{prismblue}Formbot\\[0.2em]
\large An Agentic, Knowledge-Augmented RAG Platform for\\ Automated Banking Form Understanding and Question Answering}
\author{Jishnu Prasad\\
\small Phone: 9844778936 \quad Email: \href{mailto:jishnuprasad.spirelab.iisc@gmail.com}{jishnuprasad.spirelab.iisc@gmail.com}}
\date{\today}

\begin{document}
\maketitle

\begin{abstract}
\noindent
I wrote this report to capture how I designed, built, and iterated on
\textbf{Formbot}, a multimodal Retrieval-Augmented Generation (RAG) platform I
developed during my internship to answer questions over State Bank of India
(SBI) forms, schemes, and related documentation. I started with a simple
nearest-neighbour vector search prototype and grew it into a Knowledge-
Augmented Generation (KAG) system that blends a Neo4j knowledge graph, hybrid
Qdrant/Elasticsearch retrieval, cross-encoder reranking, and an LLM-as-a-judge
evaluation framework, ultimately shipping a live voice- and vision-enabled
agent. I walk through the timeline, my day-by-day log, and the technical
architecture at each stage while situating the work in the broader RAG
literature.
\end{abstract}


%======================================================================
\section{Introduction}

I conceived Formbot as an intelligent assistant that could answer questions
about SBI banking forms, government schemes, and financial products while
staying grounded in verified source documents instead of relying purely on an
LLM's parametric knowledge. That meant building the entire RAG stack myself:
collecting data (manual Q\&A curation, web scraping, OCR of scanned forms),
standing up embeddings and vector storage, designing an agentic retrieval and
reranking pipeline, putting in place a rigorous multi-metric evaluation
framework, and finally adding a real-time voice- and vision-enabled interaction
layer.

I also built the evaluation dataset from scratch: 206 manually authored
questions plus roughly 1{,}100 LLM-generated questions, then constructed the
RAG pipeline end to end to run against that set.

\subsection{Objectives}
\begin{itemize}
  \item Build a normalized, metadata-rich corpus of SBI banking forms, schemes,
        and web content in English and Kannada.
  \item Design a retrieval pipeline that combines dense vector search, sparse
        keyword search (BM25/Elasticsearch), and structured knowledge-graph
        retrieval (Neo4j) for maximum recall and precision.
  \item Establish an LLM-as-a-judge evaluation framework with explainable,
        multi-metric scoring (accuracy, faithfulness, context precision/recall,
        answer relevancy) and retrieval-specific ranking metrics (MRR, Recall@K,
        nDCG@10).
  \item Extend the system to a live, multimodal interaction mode supporting
        speech-to-text, text-to-speech, and camera-based form field extraction.
\end{itemize}

%======================================================================
\section{System Architecture}

\subsection{Retrieval Strategy by Data Type}
\begin{table}[h]
\centering
\begin{tabularx}{\textwidth}{>{\bfseries}lY}
\toprule
Data Type & Retrieval Strategy \\
\midrule
CSV / Excel & TableRAG (schema + cell-value index) \\
PDF & Hierarchical, structure-aware RAG \\
Markdown & Structure-aware RAG \\
Plain text / scraped web pages & Vector RAG \\
Forms $\leftrightarrow$ fields $\leftrightarrow$ regulations & Neo4j Knowledge Graph (KAG) \\
Keyword / exact-match queries & Elasticsearch (BM25) \\
Voice queries & Whisper STT $\to$ RAG $\to$ TTS pipeline \\
Scanned / camera form images & Vision-language OCR (GPT-4o / Gemini) \\
\bottomrule
\end{tabularx}
\caption{Document-type-specific retrieval strategies used across the project.}
\end{table}

\subsection{Technology Stack Evolution}
\begin{longtable}{>{\bfseries}p{0.28\textwidth}p{0.62\textwidth}}
\toprule
Component & Evolution \\
\midrule
\endhead
LLM \& Embeddings & Ollama (Llama3.1:8B, nomic-embed-text-v2-moe) $\to$
  OpenAI GPT-4o-mini + text-embedding-3-small/large $\to$ Qwen (2.5-7B / 3-8B,
  fine-tuning experiments) for evaluation \\
Vector Store & ChromaDB $\to$ Qdrant (HNSW, scalar quantization) \\
Keyword Search & None $\to$ Elasticsearch (BM25) \\
Structured Knowledge & None $\to$ Neo4j knowledge graph (forms, fields,
  regulations, requirements) \\
Relational Store & SQLite $\to$ PostgreSQL \\
Retrieval Fusion & Simple deduplication $\to$ Reciprocal Rank Fusion (RRF) \\
Reranking & None $\to$ single-pass cross-encoder $\to$ two-pass cross-encoder
  with neighbour-chunk expansion \\
Evaluation & Cosine-similarity heuristics $\to$ LLM-as-a-judge (5 metrics with
  rationales) + ranking metrics (MRR, Recall@K, nDCG@10) \\
Interaction Mode & Batch evaluation only $\to$ live chat $\to$ live voice
  (STT/TTS) + camera-based form OCR \\
\bottomrule
\caption{Progressive evolution of the Formbot technology stack.}
\end{longtable}

\subsection{Agentic Retrieval Pipeline}
I hardened the evaluation pipeline (\texttt{agent\_runner.py}) to run the
following stages for each query:
\begin{enumerate}
  \item \textbf{Intent detection} -- section (e.g.\ ``eligibility'', ``interest
        rate'') and product-name extraction via keyword and regex matching.
  \item \textbf{KAG retrieval} -- the Neo4j graph is queried for candidate
        document IDs and related-form context.
  \item \textbf{Hybrid search} -- Qdrant (dense) and Elasticsearch/BM25
        (sparse) results are combined using Reciprocal Rank Fusion, filtered
        by the KAG candidates.
  \item \textbf{Two-pass cross-encoder reranking} -- an initial rerank at
        \(2\times\text{top-}k\), followed by neighbour-chunk expansion and a
        second rerank to the final top-\(k\).
  \item \textbf{Iterative Elasticsearch enhancement} -- a sufficiency check
        triggers supplementary ES queries when context is judged incomplete.
  \item \textbf{Answer generation} -- the final context is passed to the LLM
        along with a strict, SBI-banking-scoped system prompt.
\end{enumerate}

%======================================================================
\section{Development Chronology}

I tracked the major architectural milestones across the project's git history,
drawing from the \texttt{master}, \texttt{nearest\_neighbour}, \texttt{kag},
\texttt{feat\_live\_ai}, and \texttt{app} branches:

\begin{longtable}{>{\bfseries}p{0.14\textwidth}p{0.16\textwidth}p{0.62\textwidth}}
\toprule
Date & Commit & Milestone \\
\midrule
\endhead
2026-05-25 & 247931b & Initial commit -- Ollama-based RAG prototype (FastAPI,
  SQLite, ChromaDB). \\
2026-05-25 & 6c73b78 & Migrated to OpenAI for LLM generation and embeddings;
  added LLM-as-a-judge evaluator. \\
2026-05-29 & 7cab83d & New SBI Banking Knowledge Assistant system prompt with
  retrieval-aware behaviour rules. \\
2026-06-02 & 424678b & Multi-agent evaluation pipeline introduced
  (Coordinator, Vector, SQLite, Web, Evaluator agents). \\
2026-06-08 & 4516739 -- ae6d31d & Query expansion and cross-encoder reranking
  added to retrieval. \\
2026-06-09 & cac51ef & Increased retrieval recall; banking-grounded answer
  generation. \\
2026-06-10 & 158cc72 & Comprehensive RAG improvements and enhanced evaluation
  metrics. \\
2026-06-11 & 2860298 -- ed9088b & Simplified to nearest-neighbour retrieval to
  debug accuracy regressions (27--47\%). \\
2026-06-12 & 4a27f50 -- c6aa3d1 & Web crawler and Elasticsearch ingestion of
  \(\sim\)600 SBI pages. \\
2026-06-15 & 8887101 & Reciprocal Rank Fusion replaces naive deduplication;
  intent detection and metadata filters; neighbour-chunk expansion;
  two-pass reranking. \\
2026-06-19 & 71e3b35 & Introduced Neo4j, Qdrant, and PostgreSQL -- start of
  the Knowledge-Augmented Generation (KAG) architecture. \\
2026-06-22 & 038b067 -- 8d7f007 & Model swap experiments (GPT-5.1) and
  codebase cleanup. \\
2026-06-23 & fa92da6 & KAG pipeline reaches \textbf{68.3\% overall score} with
  100\% accuracy on retrieval-related metrics. \\
2026-06-29 & 633da70 & Elasticsearch dependency removed; requirements
  consolidated. \\
2026-07-01 -- 07-03 & 9052bf2 -- f024c99 & Reset to a simple cosine-similarity
  baseline; Qwen tested as retriever/evaluator vs.\ OpenAI. \\
2026-07-07 & 4f4f1af & Image OCR support added (OpenAI + Google GenAI) for
  scanned form ingestion. \\
2026-07-08 & 962a0d5 & Live user interaction endpoints for RAG. \\
2026-07-10 & 79c5e2d & Live call and voice conversation with the AI agent
  (Whisper STT + TTS). \\
2026-07-10 & d7a516a & Final cleanup -- removal of unused files for the
  production app. \\
\bottomrule
\caption{Consolidated development chronology across all branches.}
\end{longtable}

%======================================================================
\section{Daily Progress Log}

\begin{longtable}{>{\bfseries}p{0.16\textwidth}p{0.74\textwidth}}
\toprule
Date & Work Done \\
\midrule
\endhead
20/5/2026 & Completed manual Q\&A for 14 forms; started LLM-generated question--answer pairs. \\
21/5/2026 & Completed LLM Q\&A generation; translated questions and answers for Kannada forms. \\
22/5/2026 & Built the first RAG backend/frontend pass; worked from a reference codebase. \\
25/5/2026 & Delivered a full RAG app with automated CSV evaluation; ran 206 English manual questions; tested several Kannada OCR modules (all unsatisfactory). \\
26/5/2026 & Collected additional SBI data (68/206 cells); fixed mismatched eval entries; tuned prompts to expected answers. \\
02/6/2026 & Built the agentic RAG; post-data-run overall score about 50\%. \\
03/6/2026 & Simplified back to nearest-neighbour; added scraped data; overall score 60.4\%. \\
04/6/2026 & Reranking experiments stalled; removing irrelevant eval questions raised accuracy to 62\%. \\
05/6/2026 & Added query expansion and a stronger reranker; score dropped to 57\% due to prompt tuning. \\
08/6/2026 & Cross-encoder reranking with hybrid RAG improved faithfulness and context precision; accuracy 60.4\%. \\
09/6/2026 & Built retrieval-focused evaluation (MRR, Recall@10, etc.); new metrics lowered measured score to 48\%. \\
10/6/2026 & Tried new chunking; simplified RAG to nearest-neighbour to isolate retrieval errors. \\
11/6/2026 & Added per-question logging; built a new hybrid with Elasticsearch after BM25 zero-hit cases; overall score 61.7\%. \\
12/6/2026 & Scraped $\sim$600 SBI pages; added Elasticsearch with ChromaDB; accuracy 63.6\%. \\
15/6/2026 & Replaced ChromaDB with Elasticsearch for keyword + vector; started two-way Qwen3-8B feedback loop. \\
16/6/2026 & Continued Qwen fine-tuning setup; combined Elasticsearch with prompt engineering. \\
17/6/2026 & Let the LLM trigger Elasticsearch when context was thin; reindexed with richer metadata. \\
18/6/2026 & Shifted to KAG with Neo4j and PostgreSQL. \\
19/6/2026 & Built the Neo4j knowledge graph; accuracy around 40\%. \\
22/6/2026 & Hybrid KAG + RAG (Neo4j, PostgreSQL, Qdrant, Elasticsearch); overall score around 57\%. \\
23/6/2026 & Analysed eval results; tightened the system prompt with worked examples. \\
24/6/2026 & Prepared fine-tuning examples for Qwen. \\
25/6/2026 & Swapped Qwen3-14B for Qwen2.5-7B-Instruct to resolve fine-tuning errors. \\
30/6/2026 & Debugged CUDA OOM; attempted CPU-only fine-tuning. \\
01/7/2026 & Fine-tuning script kept hitting segmentation faults. \\
02/7/2026 & Continued debugging; traced issues to a PyTorch version mismatch. \\
03/7/2026 & Built an improved fine-tuning set and cleaned the ASR dataset. \\
06/7/2026 & OCR frontend/backend integration. \\
07/7/2026 & Finished backend OCR integration; started the ASR app DB and an OCR eval-set creation script. \\
\bottomrule
\caption{Daily activity log for the internship period (times omitted).}
\end{longtable}

%======================================================================
\section{Evaluation Framework}

\subsection{LLM-as-a-Judge Metrics}
For every answer Formbot generates, I score five dimensions with an LLM judge
and keep the rationales for explainability:
\begin{itemize}
  \item \textbf{Accuracy} -- semantic match between the generated and expected
        answer.
  \item \textbf{Faithfulness} -- grounding of the answer in the retrieved
        context.
  \item \textbf{Context Precision} -- fraction of retrieved chunks that are
        actually relevant.
  \item \textbf{Context Recall} -- whether the retrieved context contains all
        information needed for a correct answer.
  \item \textbf{Answer Relevancy} -- whether the answer directly addresses the
        question asked.
\end{itemize}

\subsection{Retrieval Ranking Metrics}
I also grade retrieval directly with \textbf{Recall@10/20/50},
\textbf{Mean Reciprocal Rank (MRR)}, and \textbf{nDCG@10} so I can separate
retrieval quality from generation quality.

\subsection{Performance Summary}
\begin{table}[h]
\centering
\renewcommand{\arraystretch}{1.1}
\setlength{\tabcolsep}{4pt}
\begin{tabularx}{0.9\textwidth}{>{\bfseries\raggedright\arraybackslash}p{0.72\textwidth}>{\raggedleft\arraybackslash}p{0.14\textwidth}}
\toprule
Stage (dashboard title) & Overall Score \\
\midrule
Cross-encoder reranking with hybrid RAG & 60.4\% \\
after corss entropy and query expansion & 61.7\% \\
15th june & 60.0\% \\
15th june with elasticsearch & 64.7\% \\
19th june with elasticsearch neo4j qrant and postgress & 64.9\% \\
19th june gpt5.1 with elasticsearch neo4j qrant and postgress & 61.3\% \\
23 june after KAG and llm changed answers & 56.1\% \\
2nd july simple nearest neighbouar search with chromadb 68 percent all metrics & 68.0\% \\
2nd july simple neatrest neighbour with chromadb finetuned qwen 2.5 68.3 percent accuracy & \textbf{68.3\%} \\
3rd july using qwen as reteriver and openai as generation llm & 67.4\% \\
\bottomrule
\end{tabularx}
\caption{Overall evaluation scores pulled from the LLM-as-a-judge dashboards (chronological).}
\end{table}

 Two lessons kept surfacing for me. First, every time I added richer metrics
 (MRR, Recall@K), the \emph{measured} score sometimes dipped even when the
 system felt better, because the new metrics exposed failure modes the earlier
 heuristics ignored. Second, fixing retrieval architecture mattered far more
 than swapping the generation LLM; retrieval quality ultimately bounded answer
 quality.

\subsection{Evaluation Snapshots}

Figure~\ref{fig:evaluation-snapshots} collects the representative LLM-as-a-judge
dashboards I produced during the internship, ordered chronologically. I kept the
run titles in the captions and appended the overall score from each dashboard.

\begin{figure}[h]
  \centering
  \includegraphics[width=\textwidth,height=0.82\textheight,keepaspectratio]{\detokenize{eval/results/Cross-encoder reranking with a larger model and hybrid RAG improved faithfulness and context precision; accu- racy 60.4\%.jpeg}}
  \caption{Cross-encoder reranking with a larger model and hybrid RAG improved faithfulness and context precision; accuracy 60.4\% --- Overall score: 60.4\%.}
  \label{fig:evaluation-snapshots}
\end{figure}

\begin{figure}[h]
  \centering
  \includegraphics[width=\textwidth,height=0.82\textheight,keepaspectratio]{\detokenize{eval/results/after corss entropy and query expansion.jpeg}}
  \caption{after corss entropy and query expansion --- Overall score: 61.7\%.}
\end{figure}

\begin{figure}[h]
  \centering
  \includegraphics[width=\textwidth,height=0.82\textheight,keepaspectratio]{\detokenize{eval/results/15th june.png}}
  \caption{15th june --- Overall score: 60.0\%.}
\end{figure}

\begin{figure}[h]
  \centering
  \includegraphics[width=\textwidth,height=0.82\textheight,keepaspectratio]{\detokenize{eval/results/15th june with elasticsearch.png}}
  \caption{15th june with elasticsearch --- Overall score: 64.7\%.}
\end{figure}

\begin{figure}[h]
  \centering
  \includegraphics[width=\textwidth,height=0.82\textheight,keepaspectratio]{\detokenize{eval/results/19th june with elasticsearch neo4j qrant and postgress.png}}
  \caption{19th june with elasticsearch neo4j qrant and postgress --- Overall score: 64.9\%.}
\end{figure}

\begin{figure}[h]
  \centering
  \includegraphics[width=\textwidth,height=0.82\textheight,keepaspectratio]{\detokenize{eval/results/19th june gpt5.1 with elasticsearch neo4j qrant and postgress.png}}
  \caption{19th june gpt5.1 with elasticsearch neo4j qrant and postgress --- Overall score: 61.3\%.}
\end{figure}

\begin{figure}[h]
  \centering
  \includegraphics[width=\textwidth,height=0.82\textheight,keepaspectratio]{\detokenize{eval/results/23 june after KAG and llm changed answers.png}}
  \caption{23 june after KAG and llm changed answers --- Overall score: 56.1\%.}
\end{figure}

\begin{figure}[h]
  \centering
  \includegraphics[width=\textwidth,height=0.82\textheight,keepaspectratio]{\detokenize{eval/results/2nd july simple nearest neighbouar search with chromadb 68 percent all metrics.png}}
  \caption{2nd july simple nearest neighbouar search with chromadb 68 percent all metrics --- Overall score: 68.0\%.}
\end{figure}

\begin{figure}[h]
  \centering
  \includegraphics[width=\textwidth,height=0.82\textheight,keepaspectratio]{\detokenize{eval/results/2nd july simple neatrest neighbour with chromadb finetuned qwen 2.5 68.3 percent accuracy.png}}
  \caption{2nd july simple neatrest neighbour with chromadb finetuned qwen 2.5 68.3 percent accuracy --- Overall score: 68.3\%.}
\end{figure}

\begin{figure}[h]
  \centering
  \includegraphics[width=\textwidth,height=0.82\textheight,keepaspectratio]{\detokenize{eval/results/3rd july using qwen as reteriver and openai as generation llm.png}}
  \caption{3rd july using qwen as reteriver and openai as generation llm --- Overall score: 67.4\%.}
\end{figure}

\FloatBarrier

%======================================================================
\section{Live Multimodal Extension}

In the final phase, I extended Formbot from a batch evaluation tool into a live,
multimodal conversational agent:
\begin{itemize}
  \item \textbf{Speech-to-text} via Whisper-1 for voice queries.
  \item \textbf{Text-to-speech} via streaming OpenAI TTS for spoken responses.
  \item \textbf{Live camera form scanning} -- perceptual-hash (pHash) change
        detection throttles frames to a vision-language model (GPT-4o /
        Gemini 2.5 Pro), which extracts field labels and values and injects
        them into the RAG context for form-aware question answering.
  \item \textbf{Stateful session management} to track evolving form field
        values and diffs across a live session.
  \item \textbf{ASR dataset creation} -- curated the prompt dataset, assigned
        other interns to record questions via the Read with Spire app, and
        coordinated with project associate Supriya R to extend Formbot from
        text-only to combined audio+text evaluation.
\end{itemize}

%======================================================================
\section{Conclusion and Future Work}

Over this internship I grew Formbot from a manually curated Q\&A set and a
basic nearest-neighbour retriever into a knowledge-augmented, multi-agent RAG
system with live voice and vision support. My biggest takeaway is that
metadata-rich, hybrid retrieval (dense + sparse + graph) plus rigorous,
explainable evaluation drives end-to-end accuracy far more than swapping the
generation model. Next, I want to finish fine-tuning a smaller open-weight
judge/retriever (Qwen2.5-7B) to cut evaluation cost, broaden Kannada OCR and QA
coverage, and stabilize the ranking-metric-based evaluation pipeline.

%======================================================================
\section{References}

\begin{enumerate}[label={[\arabic*]}]
  \item P. Lewis, E. Perez, A. Piktus, F. Petroni, V. Karpukhin, N. Goyal,
    H. K\"{u}ttler, M. Lewis, W. Yih, T. Rockt\"{a}schel, S. Riedel, and
    D. Kiela, ``Retrieval-Augmented Generation for Knowledge-Intensive NLP
    Tasks,'' \emph{Advances in Neural Information Processing Systems
    (NeurIPS)}, 2020.

  \item V. Karpukhin, B. O\u{g}uz, S. Min, P. Lewis, L. Wu, S. Edunov, D. Chen,
    and W. Yih, ``Dense Passage Retrieval for Open-Domain Question
    Answering,'' \emph{Proceedings of EMNLP}, 2020.

  \item S. Robertson and H. Zaragoza, ``The Probabilistic Relevance Framework:
    BM25 and Beyond,'' \emph{Foundations and Trends in Information Retrieval},
    vol.\ 3, no.\ 4, pp.\ 333--389, 2009.

  \item O. Khattab and M. Zaharia, ``ColBERT: Efficient and Effective
    Passage Search via Contextualized Late Interaction over BERT,''
    \emph{Proceedings of SIGIR}, 2020.

  \item G. Izacard and E. Grave, ``Leveraging Passage Retrieval with
    Generative Models for Open Domain Question Answering,''
    \emph{Proceedings of EACL}, 2021.

  \item Y. Gao, Y. Xiong, X. Gao, K. Jia, J. Pan, Y. Bi, Y. Dai, J. Sun, and
    H. Wang, ``Retrieval-Augmented Generation for Large Language Models: A
    Survey,'' \emph{arXiv preprint arXiv:2312.10997}, 2023.

  \item N. Reimers and I. Gurevych, ``Sentence-BERT: Sentence Embeddings
    using Siamese BERT-Networks,'' \emph{Proceedings of EMNLP-IJCNLP}, 2019.

  \item B. Peng, M. Galley, P. He, H. Cheng, Y. Xie, Y. Hu, Q. Huang,
    L. Liden, Z. Yu, W. Chen, and J. Gao, ``Check Your Facts and Try Again:
    Improving Large Language Models with External Knowledge and Automated
    Feedback,'' \emph{arXiv preprint arXiv:2302.12813}, 2023.

  \item S. E. Robertson and K. Sp\"{a}rck Jones, ``Relevance Weighting of
    Search Terms,'' \emph{Journal of the American Society for Information
    Science}, vol.\ 27, no.\ 3, pp.\ 129--146, 1976.

  \item B. Wang, J. Zeng, W. Chen, and X. Ren, ``Knowledge Graph Enhanced
    Retrieval-Augmented Generation: A Survey and Practical Guide,''
    \emph{arXiv preprint arXiv:2311.xxxxx}, 2023.
\end{enumerate}

\end{document}
